\section{Enveloping algebras and primitive cohomology}
\label{sp:section}

Let $\kk$ be a field of characteristic zero. For a coproduct with
distinguished element $1$, an element $z$ is primitive when
$\Delta z=z\otimes1+1\otimes z$.
A Hall product is an associative operation. A primitive Lie algebra
also requires a coproduct. Even the algebra $\kk[x,y]$ permits different
answers. Give $x$ height one and $y$ height two. One coproduct makes
both generators primitive; another makes $x$ primitive and sets
\[
 \Delta_1y=y\otimes1+1\otimes y+x\otimes x.
\]
Both extend multiplicatively, with counits killing $x,y$.
The two iterated coproducts of $y$ have the same three cross terms,
so the second coproduct is coassociative. In it, $y$ is not primitive
but $y-\tfrac12x^2$ is. Thus the product alone does not select the
primitive subspace.

An enveloping algebra has a specified coproduct. We prove the ordered
basis theorem, identify its primitives, and compare its cohomology with
the enveloping algebra of the Lie cohomology. These results will supply
the algebraic comparison used in positive Borcherds recognition and in
the current examples. Their hypotheses do not assert a geometric Hall
comparison.

\subsection{Odd squares and ordered products}

To retain signs in an associative multiplication, distinguish even and
odd vectors.
A super vector space is a vector space with a decomposition into even
and odd parts.  The parity of a homogeneous vector $x$ is
$p(x)\in\mathbb Z/2$.  Interchanging homogeneous factors uses the
Koszul sign $(-1)^{p(x)p(y)}$.
An associative superalgebra has a unit of even parity and a multiplication
that adds parities.  Its graded commutator is
\[
 [x,y]_{\mathrm{ass}}=xy-(-1)^{p(x)p(y)}yx.
\]
For an odd element $x$ this gives $[x,x]_{\mathrm{ass}}=2x^2$.
Its square can therefore be nonzero.  Forming an associative algebra
from a graded bracket must retain this odd self-bracket relation.


For a first ordered basis, take a generator $e$ of degree zero and a generator
$\theta$ of degree one, with their parities given by degree modulo two.
Impose the relations $e\theta=\theta e$ and $\theta^2=0$.
Every word reduces to $e^m$ or $e^m\theta$, or to zero. These expressions
are independent: on $\kk[e]\oplus\kk[e]\theta$ the multiplication
\[
 (a+b\theta)(c+f\theta)=ac+(af+bc)\theta
\]
is associative and has those relations. The maps from the word quotient
to this algebra and back fix both generators, hence are inverse.
This constructs the basis of $\kk[e]\otimes\Lambda(\theta)$.
\label{sp:even-odd-model}

A Lie superalgebra $\mathfrak g$ has an even bilinear bracket satisfying
\begin{align}
 [x,y]&=-(-1)^{p(x)p(y)}[y,x],\\
 [x,[y,z]]&=[[x,y],z]+(-1)^{p(x)p(y)}[y,[x,z]].
 \label{sp:jacobi}
\end{align}
These identities apply to homogeneous vectors and extend by linearity.
For an odd vector $x$, they imply
\begin{equation}
 2[x,[x,z]]=[[x,x],z],\qquad [x,[x,x]]=0.
 \label{sp:odd-jacobi}
\end{equation}
The first identity follows by setting $y=x$ in \eqref{sp:jacobi}.
For $z=x$, skew symmetry turns it into $3[x,[x,x]]=0$.
This proves the second identity.

The tensor algebra
$T(\mathfrak g)=\bigoplus_{r\geq0}\mathfrak g^{\otimes r}$
has concatenation as multiplication and $\mathfrak g^{\otimes0}=\kk1$.
Every element is a finite sum of finite words.  The enveloping algebra is
\begin{equation}
 U(\mathfrak g)=T(\mathfrak g)/(R(x,y)),\qquad
 R(x,y)=xy-(-1)^{p(x)p(y)}yx-[x,y].
 \label{sp:relations}
\end{equation}
The parentheses denote the two-sided ideal generated by all the
displayed relations.  A parity-preserving linear map from $\mathfrak g$
to an associative superalgebra extends uniquely to $T(\mathfrak g)$.
It factors through $U(\mathfrak g)$ precisely when it preserves brackets,
using the graded commutator in the target.  This proves the universal
property directly from the quotient.

\begin{example}[An odd generator with a square]\label{sp:odd-square}
Take an odd basis vector $x$ and an even basis vector $t$, with
$[x,x]=2t$ and $t$ central.  These formulas satisfy
\eqref{sp:jacobi}, since every nested bracket vanishes.
Then
\[
 U(\mathfrak g)=\kk\langle t,x\rangle/(tx-xt,\ x^2-t).
\]
The homomorphism to the associative polynomial algebra $\kk[X]$ defined
by $x\mapsto X$, $t\mapsto X^2$ has inverse $X\mapsto x$.
Here $X$ has odd parity.  This algebra is not supercommutative.
The vectors $t^m$ and $t^mx$, $m\geq0$, form a basis.
In contrast, for one odd vector with zero bracket the relation is
$x^2=0$, and the enveloping algebra is $\kk\oplus \kk x$.
\end{example}

Give $U(\mathfrak g)$ the increasing word filtration $F_r$, consisting
of images of words of length at most $r$, with $F_{-1}=0$.
Its associated graded vector space is
$\operatorname{gr}_F U=\bigoplus_{r\geq0}F_r/F_{r-1}$.
Since $F_rF_s\subseteq F_{r+s}$, multiplication induces a product
$(F_r/F_{r-1})\otimes(F_s/F_{s-1})\to F_{r+s}/F_{r+s-1}$.
Changing a representative adds a term in $F_{r+s-1}$, so this product
is well-defined. It makes the associated graded space an algebra.
Let
\[
 S(\mathfrak g)=T(\mathfrak g)/(xy-(-1)^{p(x)p(y)}yx)
             =\bigoplus_{r\geq0}S^r(\mathfrak g)
\]
be the super symmetric algebra.  An odd variable has square zero here.
The highest-length parts of \eqref{sp:relations} give a surjective map
$S(\mathfrak g)\to\operatorname{gr}_F U(\mathfrak g)$.

\begin{theorem}[The super Poincar\'e--Birkhoff--Witt theorem]
\label{sp:pbw}
Choose a totally ordered homogeneous vector-space basis $(x_i)_{i\in I}$
of $\mathfrak g$.  The words
\[
 x_{i_1}\cdots x_{i_r},\qquad i_1\leq\cdots\leq i_r,
 \qquad i_j<i_{j+1}\text{ if }x_{i_j}\text{ is odd},
\]
together with the empty word, form a basis of $U(\mathfrak g)$.
The map $\mathfrak g\to U(\mathfrak g)$ is injective, and
$S(\mathfrak g)\to\operatorname{gr}_F U(\mathfrak g)$ is an
isomorphism.  No dimension bound on $\mathfrak g$ is required.
\end{theorem}

\begin{proof}
Replace an adjacent inverted pair by
\[
 x_jx_i\longmapsto(-1)^{p(x_j)p(x_i)}x_ix_j+[x_j,x_i]
 \quad(j>i).
\]
Replace an adjacent equal odd pair by
$xx\longmapsto\tfrac12[x,x]$.
Expand each bracket in the chosen basis.  Each resulting term has
smaller word length, or the same length and fewer inverted pairs.
Ordering these two nonnegative integers lexicographically proves
termination.  Each expansion is finite, including for an infinite basis.

Independence of the replacements requires checking overlapping pairs.
Replacements on disjoint pairs commute.  Both adjacent pairs in a
word of length three can be replaceable in exactly four ways.

First take distinct basis vectors $a>b>c$.  Write
$A=(-1)^{p(a)p(b)}$, $B=(-1)^{p(a)p(c)}$, and
$C=(-1)^{p(b)p(c)}$.  Starting at the left or right pair, and moving
the original three factors into order, gives respectively
\begin{align*}
 ABCcba+AB[b,c]a+A b[a,c]+[a,b]c,\\
 ABCcba+BCc[a,b]+C[a,c]b+a[b,c].
\end{align*}
Their difference has length at most two.  Using the relations on two inputs, this difference becomes
\[
 -[a,[b,c]]+A[b,[a,c]]+[[a,b],c]=0.
\]
The equality is \eqref{sp:jacobi}.

Next consider $xxc$, with $x$ odd and $x>c$.  Put
$\epsilon=(-1)^{p(c)}$.  The two expressions after the initial
replacements are
\[
 \tfrac12[x,x]c,\qquad
 \tfrac12c[x,x]+\epsilon[x,c]x+x[x,c].
\]
Their difference reduces to
$\tfrac12[[x,x],c]-[x,[x,c]]=0$ by \eqref{sp:odd-jacobi}.
For $axx$, with $a>x$ and $x$ odd, put
$\epsilon=(-1)^{p(a)}$.  The two expressions are
\[
 \tfrac12[x,x]a+\epsilon x[a,x]+[a,x]x,\qquad
 \tfrac12a[x,x].
\]
Their difference reduces to
$[[a,x],x]-\tfrac12[a,[x,x]]=0$.
Indeed \eqref{sp:jacobi} and skew symmetry give
$[a,[x,x]]=2[[a,x],x]$.
Finally, for three equal odd factors $xxx$, the difference is
$\tfrac12[[x,x],x]=0$ by \eqref{sp:odd-jacobi}.

The relations on two inputs used here hold under normalization for any
homogeneous linear combinations.  For an increasing pair, this follows
from its reversed pair and skew symmetry.  For an equal even pair,
both the commutator and bracket vanish.  For an equal odd pair, it is
exactly the half-bracket rule.

For completeness, these overlap calculations imply uniqueness by an
induction on length and inversions.  If two first replacements differ,
disjoint ones commute.  Overlapping ones give the differences just
calculated.  Their remaining terms have strictly smaller length, so the
induction hypothesis identifies their normalized expressions.  The same
argument applies inside any fixed left and right word context.
All bracket corrections then have smaller total length, and each
adjacent swap has fewer inversions.  This proves uniqueness for the
whole word and for its linear combinations.

Thus normalization is a linear map to the span of the stated words.
It annihilates each defining relation in every word context, and hence
the ideal in \eqref{sp:relations}.  It fixes every ordered word.
These words are therefore linearly independent in the quotient, as
well as spanning it.  Setting every bracket to zero in the same
argument gives the corresponding basis of $S(\mathfrak g)$.
The bases in each length identify $S(\mathfrak g)$ with
$\operatorname{gr}_F U(\mathfrak g)$.
\end{proof}

In Example~\ref{sp:odd-square}, $x^2=t$ lies in $F_1$ although its
displayed expression has length two.  Its degree-two symbol is zero.
This explains the exterior odd factor in $S(\mathfrak g)$ without
discarding the nonzero element $t$ of $U(\mathfrak g)$.

\subsection{The differential and symmetrization}
\label{sp:differential-section}

A cochain complex $V=\bigoplus_{n\in\mathbb Z}V^n$ has a linear map
$d:V^n\to V^{n+1}$ with $d^2=0$.
Its parity is $p(v)=|v|\bmod2$, where $|v|=n$ for $v\in V^n$.
A differential graded Lie algebra is such a complex with a
degree-zero Lie superbracket and
\[
 d[x,y]=[dx,y]+(-1)^{|x|}[x,dy].
\]
All tensor products below are algebraic.  Their degree-$n$ part is the
direct sum over the degrees of the factors that add to $n$.
For two factors the differential is
$d(x\otimes y)=dx\otimes y+(-1)^{|x|}x\otimes dy$.
Complexes may be unbounded in both directions.

Extend $d$ to $T(\mathfrak g)$ by the graded Leibniz rule.
Its square is zero on generators, and the two mixed terms in $d^2(ab)$
cancel, so its square is zero everywhere.
Direct expansion of \eqref{sp:relations} gives
\[
 dR(x,y)=R(dx,y)+(-1)^{|x|}R(x,dy).
\]
Thus $d$ descends to $U(\mathfrak g)$ and preserves the word filtration.
The corresponding extension to $S(\mathfrak g)$ preserves symmetric
length.  Theorem~\ref{sp:pbw} identifies the associated graded
differential with this differential on $S(\mathfrak g)$.

To identify the Lie generators inside the enveloping algebra, we
also use a coproduct. A counital coaugmented coalgebra is a vector
space $C$ with linear maps $\Delta:C\to C\otimes_\kk C$ and
$\varepsilon:C\to \kk$, and a distinguished element $1$ with
$\Delta1=1\otimes1$ and $\varepsilon(1)=1$.
The coproduct and counit satisfy
\[
 (\Delta\otimes1)\Delta=(1\otimes\Delta)\Delta,\qquad
 (\varepsilon\otimes1)\Delta=(1\otimes\varepsilon)\Delta=1_C.
\]
The last two equalities use $\kk\otimes C=C=C\otimes \kk$.
Its primitive subspace is
\[
 P(C)=\{z\in C:\Delta z=z\otimes1+1\otimes z\}.
\]

The counit equations imply $\varepsilon(z)=0$ for a primitive $z$.
On $\ker\varepsilon$ the reduced coproduct is
$\bar\Delta z=\Delta z-z\otimes1-1\otimes z$. Its kernel is $P(C)$.
Here primitive describes a coproduct. A primitive for a differential,
such as an element $h$ satisfying $dh=a$, is a different use of the word.

A differential graded bialgebra is also an associative algebra whose
coproduct and counit are algebra chain maps.  Tensor multiplication is
\[
 (a\otimes b)(c\otimes e)=(-1)^{|b||c|}ac\otimes be.
\]
Associativity gives
\[
 [x,yz]_{\mathrm{ass}}=[x,y]_{\mathrm{ass}}z
       +(-1)^{|x||y|}y[x,z]_{\mathrm{ass}}.
\]
Expanding the two products proves this formula.  Substituting the
graded commutator for $yz$ proves \eqref{sp:jacobi} for this bracket.

For two homogeneous primitives, the mixed terms in the coproduct of
$xy-(-1)^{|x||y|}yx$ cancel.  Thus the primitives form a differential
graded Lie algebra.  This also proves closure under $d$ when $\Delta$
is a chain map.

On the tensor algebra define $\Delta x=x\otimes1+1\otimes x$ and
$\varepsilon(x)=0$, and extend these maps multiplicatively.
One obtains
\[
 \Delta R(x,y)=R(x,y)\otimes1+1\otimes R(x,y),
 \qquad \varepsilon R(x,y)=0.
\]
The maps therefore descend to $U(\mathfrak g)$.
Coassociativity, counitality and compatibility with $d$ hold on
generators, hence on every product.  The same construction applies to
$S(\mathfrak g)$.  Both coproducts are graded cocommutative, meaning
that the Koszul interchange of the two output factors fixes them.

These bialgebras have antipodes.  Here is a construction that requires
no finiteness of $\mathfrak g$.  Put $U=U(\mathfrak g)$.
For degree-zero linear maps $f,g:U\to U$, define
their convolution by $f*g=m(f\otimes g)\Delta$, where $m$ is
multiplication.  Coassociativity of $\Delta$ and associativity of $m$
give $(f*g)*h=f*(g*h)$ by expansion.
Its identity is $u=\eta\varepsilon$, with
$\eta:\kk\to U$ the unit map. Put $p=1_U-u$ and set $p^{*0}=u$.  On a word of length at
most $r$, $p^{*n}=0$ for $n>r$: each surviving coproduct factor must
contain a generator.  Therefore
\[
 S=\sum_{n\geq0}(-1)^np^{*n}
\]
is a well-defined linear map.  Multiplying this finite geometric series
by $u+p=1_U$ in either convolution order gives $S*1_U=1_U*S=u$.
These are the two antipode identities.  Every summand is a chain map,
so $S$ commutes with $d$.  A bialgebra with such a map is a Hopf algebra.

The filtration gives the ordered basis, but a chain isomorphism needs
a splitting that respects the differential.  Graded averaging supplies
one without requiring the chosen basis to consist of cycles.
For homogeneous $x_1,\ldots,x_r$, let
$\epsilon(\pi;x)$ be the product of Koszul signs obtained when a
permutation $\pi$ reorders these labeled factors as
$x_{\pi(1)},\ldots,x_{\pi(r)}$.
Adjacent interchanges satisfy their defining permutation relations:
two equal interchanges give sign one, disjoint interchanges commute,
and both three-factor interchanges have sign
$(-1)^{|x_1||x_2|+|x_1||x_3|+|x_2||x_3|}$.
Thus the sign does not depend on the sequence of adjacent interchanges.
Define symmetrization by
\begin{equation}
 \sigma_r(x_1\cdots x_r)=\frac1{r!}
 \sum_{\pi\in\mathfrak S_r}\epsilon(\pi;x)
 x_{\pi(1)}\cdots x_{\pi(r)},\qquad \sigma_0(1)=1.
 \label{sp:symmetrization}
\end{equation}
Here $\mathfrak S_r$ is the group of permutations of $r$ positions.
The input product belongs to $S^r(\mathfrak g)$ and the output to
$U(\mathfrak g)$.  Reindexing the sum proves compatibility with the
supercommutative input relations, so this is a well-defined linear map.

\begin{proposition}[Symmetrization as a coalgebra chain map]
\label{sp:symmetrization-theorem}
The sum $\sigma=\bigoplus_r\sigma_r:S(\mathfrak g)\to U(\mathfrak g)$
is a natural filtered isomorphism of differential graded coalgebras.
Its associated graded map is the identity under Theorem~\ref{sp:pbw}.
It is not asserted to preserve multiplication.
\end{proposition}

\begin{proof}
In the associated graded algebra every summand of
\eqref{sp:symmetrization} equals its input monomial: the sign for
permutation and the sign for reordering back cancel.
Thus the associated graded map is the identity.  A filtered element
has finite word length.  Subtracting the image of its highest symbol
and proceeding by decreasing length proves surjectivity.
The highest nonzero symbol of an element in the kernel would map to
itself, proving injectivity.  This argument uses neither completeness
nor finite-dimensional filtered pieces.

For a homogeneous pair of degrees $p,q$, the Koszul interchange
$\tau(x\otimes y)=(-1)^{pq}y\otimes x$ is a chain map, since both
$d\tau(x\otimes y)$ and $\tau d(x\otimes y)$ equal
\[
 (-1)^{pq}dy\otimes x+(-1)^{pq+q}y\otimes dx.
\]
Every permutation is consequently a chain map.  Multiplication is a
chain map by the Leibniz rule.  Formula~\eqref{sp:symmetrization}
therefore commutes with $d$.

For a subset $J\subseteq\{1,\ldots,r\}$, list the entries of $J$
and its complement in their original orders.  Let $\epsilon(J,J^c;x)$
be the sign that moves the entries of $J$ before those of $J^c$.
The primitive coproduct on the symmetric algebra is
\[
 \Delta(x_1\cdots x_r)=
 \sum_J\epsilon(J,J^c;x)x_J\otimes x_{J^c}.
\]
In $\Delta\sigma_r$, fix $J$ of size $a$ and fix orders on its two
parts.  There are $\binom ra$ interleavings giving these orders.
Their coefficient after averaging is
$\binom ra/r!=1/(a!(r-a)!)$.
Their permutation and coproduct signs combine to the subset sign
times the two internal permutation signs.  These are exactly the
terms in $(\sigma\otimes\sigma)\Delta$.
This proves coproduct compatibility, including all odd signs.
Unit and counit compatibility follow in length zero.
Finally the same formula commutes with every differential graded Lie
map, proving naturality.
\end{proof}

\begin{theorem}[Primitives of an enveloping algebra]
\label{sp:primitives}
For every differential graded Lie algebra over $\kk$,
\[
 P(U(\mathfrak g))=\mathfrak g
\]
through the canonical inclusion.  This is an equality of differential
graded Lie algebras.  With the differential omitted, the same statement
holds for Lie superalgebras.
\end{theorem}

\begin{proof}
On symmetric length $r>0$, take the coproduct component
\[
 \delta_{1,r-1}:S^r(\mathfrak g)
 \longrightarrow \mathfrak g\otimes S^{r-1}(\mathfrak g).
\]
For a monomial, each of its $r$ labeled factors can occupy the first
output position.  Multiplying the outputs restores the original
monomial: the sign moving the factor out cancels the sign moving it
back.  Therefore $m\delta_{1,r-1}=r\,1_{S^r(\mathfrak g)}$.
This identity also holds for zero monomials with repeated odd factors.
Since $r$ is invertible in $\kk$, this coproduct component is injective.
For $r\geq2$ it is a component of the reduced coproduct, so there is
no primitive element of that length.  A primitive scalar is zero by
the counit.  Hence $P(S(\mathfrak g))=\mathfrak g$.
The coalgebra isomorphism of
Proposition~\ref{sp:symmetrization-theorem} fixes $\mathfrak g$, giving
the stated equality.  The bracket and differential restrict to their
original values by construction.
\end{proof}

\subsection{Cohomology of the enveloping algebra}

For a cochain complex $V$, write
$H^n(V)=\ker(d:V^n\to V^{n+1})/\operatorname{im}(d:V^{n-1}\to V^n)$
and $H(V)=\bigoplus_nH^n(V)$.
For a differential graded Lie algebra the bracket of cycle
representatives defines a bracket on $H(\mathfrak g)$.
It is well-defined because $[da,y]=d[a,y]$ when $dy=0$.
Thus $H(\mathfrak g)$ is a graded Lie algebra with zero differential.
Multiplication likewise descends to $H(U(\mathfrak g))$.
The cycle inclusion gives an algebra map
\begin{equation}
 \Phi_{\mathfrak g}:U(H(\mathfrak g))\longrightarrow H(U(\mathfrak g)),
 \qquad [x_1]\cdots[x_r]\longmapsto[x_1\cdots x_r].
 \label{sp:cohomology-map}
\end{equation}
The enveloping relations descend because the commutator of two cycle
classes is the class of their Lie bracket.  Replacing any cycle by a
boundary changes its product with the other cycles by a boundary.
This proves the map is defined before any isomorphism is claimed.

\begin{lemma}[Tensor and symmetric powers of complexes]
\label{sp:kunnethtensor}
For complexes of vector spaces, the cross product of cycles gives
natural isomorphisms
\[
 H(V)\otimes H(W)\xrightarrow{\sim}H(V\otimes W),\qquad
 S^r(H(V))\xrightarrow{\sim}H(S^r(V)).
\]
The tensor product on each side uses direct sums of cohomological
degrees.  No boundedness assumption is needed.  The second assertion
uses characteristic zero.
\end{lemma}

\begin{proof}
In each degree, choose complements to the boundaries in the cycles,
and to the cycles in the whole space.  This gives
$V=H_V\oplus K_V$ as complexes, with $d=0$ on $H_V$.
The other summand has the form $K_V=B\oplus E$, with
$d:E^n\to B^{n+1}$ an isomorphism.  Define a degree-$-1$ map $h$
as its inverse on $B$ and zero on $E$.  Then $dh+hd=1$ on $K_V$.
On $K_V\otimes W$, the map $h\otimes1$ is a contraction: the terms
involving $d_W$ cancel because $h$ is odd.
On $W\otimes K_V$, the contraction is
$w\otimes v\mapsto(-1)^{|w|}w\otimes h(v)$.
Splitting both complexes thus leaves only $H_V\otimes H_W$ in
cohomology.  The resulting isomorphism is the stated cross product,
so the map is independent of the complements.
The same argument applies to any finite tensor power and gives
$H(V)^{\otimes r}\cong H(V^{\otimes r})$.
Every vector has a finite expansion in the direct-sum tensor product,
so these contractions are valid for unbounded complexes.

On $V^{\otimes r}$ the signed permutations are chain maps, as checked
in Proposition~\ref{sp:symmetrization-theorem}.
The averaging operator $e_r=(1/r!)\sum_{\pi\in\mathfrak S_r}\pi$
is a chain map with $e_r^2=e_r$, hence an idempotent.
Its kernel is the span of $v-\pi v$:
these differences are in the kernel, and
$v-e_rv=(1/r!)\sum_\pi(v-\pi v)$ proves the reverse inclusion.
Consequently $S^r(V)$, the quotient by those differences, identifies
with $\operatorname{im}e_r$ as a complex.
An idempotent splits a complex into its image and kernel, so
$H(\operatorname{im}e_r)=\operatorname{im}H(e_r)$.
The tensor cross product commutes with signed permutations of cycles.
Taking the image of this idempotent proves the second isomorphism.
\end{proof}

This first isomorphism is the algebraic K\"unneth isomorphism.
Applied to the coproduct, it defines the graded Hopf structure on
$H(U(\mathfrak g))$.  Its structure equations follow by applying
cohomology and the natural cross products to the original equations.

\begin{theorem}[Enveloping algebras and primitive cohomology]
\label{sp:cohomology}
For every differential graded Lie algebra $\mathfrak g$ over $\kk$,
the map \eqref{sp:cohomology-map} is a natural graded Hopf isomorphism.
The inclusion of strict primitives induces a natural isomorphism
\begin{equation}
 H(P(U(\mathfrak g)))\xrightarrow{\sim}
 P(H(U(\mathfrak g)))\cong H(\mathfrak g).
 \label{sp:primitive-cohomology}
\end{equation}
The first map sends the class of a primitive cycle to its class in
the whole enveloping algebra.  It preserves the graded Lie bracket.
\end{theorem}

\begin{proof}
Cohomology commutes with direct sums of complexes: both cycles and
boundaries are computed summand by summand.  The preceding lemma
therefore gives a natural isomorphism
$S(H(\mathfrak g))\to H(S(\mathfrak g))$.
There is a commutative diagram of graded vector spaces
\[
 \begin{array}{ccc}
 S(H(\mathfrak g))&\xrightarrow{\ \sim\ }&H(S(\mathfrak g))\\
 {\scriptstyle\sigma}\big\downarrow&&
     \big\downarrow{\scriptstyle H(\sigma)}\\
 U(H(\mathfrak g))&\xrightarrow{\ \Phi_{\mathfrak g}\ }&H(U(\mathfrak g)).
 \end{array}
\]
On each product of cycle classes, both paths take the same signed
average of products of representatives.  This proves commutativity.
The other three arrows are isomorphisms by
Proposition~\ref{sp:symmetrization-theorem} and
Lemma~\ref{sp:kunnethtensor}, so $\Phi_{\mathfrak g}$ is an isomorphism.
It is an algebra map by construction.  It preserves coproducts since
the generating classes $[x]$ are primitive on both sides and the
coproducts are algebra maps.  It also preserves counits and antipodes.
For antipodes this follows from uniqueness of a convolution inverse.

Theorem~\ref{sp:primitives} identifies the source of
\eqref{sp:primitive-cohomology} with $H(\mathfrak g)$.
Applying the same theorem to the zero-differential graded Lie algebra
$H(\mathfrak g)$ and using $\Phi_{\mathfrak g}$ identifies the target
with $H(\mathfrak g)$.  Under these identifications the first map is
the identity, since every cycle $x$ maps to its own class $[x]$.
Naturality and bracket compatibility follow on these representatives.
\end{proof}

\begin{example}[A differential on an odd square]\label{sp:odd-differential}
In Example~\ref{sp:odd-square}, give $x$ degree one, $t$ degree two,
and set $dx=t$, $dt=0$.  The Lie derivation rule on $[x,x]=2t$ gives
$[t,x]-[x,t]=0$, so this is a differential graded Lie algebra.
Its complex $\kk x\xrightarrow{d}\kk t$ is acyclic.
In its enveloping algebra $t=x^2$ and
\[
 d(x^n)=
 \begin{cases}x^{n+1},&n\text{ odd},\\0,&n\text{ even}.
 \end{cases}
\]
The formula follows by summing the $n$ alternating Leibniz terms.
All positive-degree cohomology vanishes and $H(U(\mathfrak g))=\kk$,
as \eqref{sp:cohomology-map} predicts.
The coproduct gives
$\Delta(x^2)=x^2\otimes1+1\otimes x^2$, because its mixed terms cancel.
Thus the strict primitive complex is exactly
$\kk x\xrightarrow{d}\kk x^2$, and both sides of
\eqref{sp:primitive-cohomology} vanish.
No assignment of positive integer weights to $x$ and $t$ makes $d$
preserve weight and the bracket add weights. The equation $dx=t$
requires equal weights. The equation $[x,x]=2t$ requires the weight of
$t$ to be twice that of $x$. These requirements are incompatible with
a positive weight for $x$.
\end{example}

\begin{example}[An abelian odd class and a contractible pair]
For $\mathfrak g=\kk x$ with $|x|=1$, zero bracket and zero differential,
$U(\mathfrak g)=\kk\oplus \kk x$ and $P(U(\mathfrak g))=\kk x$.
The cohomology comparison fixes this odd class.
For an abelian complex $\kk u\xrightarrow{d}\kk v$, with $|u|=0$,
$|v|=1$ and $du=v$, one has
\[
 U(\mathfrak g)=\kk[u]\otimes\Lambda(v),\qquad
 d(u^m)=m u^{m-1}v,\qquad d(u^mv)=0.
\]
Here $\Lambda(v)=\kk\oplus \kk v$ with $v^2=0$.
Every $u^mv$ is the boundary of $u^{m+1}/(m+1)$, and a polynomial
of degree zero in cohomology is closed precisely when it is constant.
Thus again $H(U(\mathfrak g))=\kk$.
\end{example}

\subsection{Lie quotients}

Let $\mathfrak r\subseteq\mathfrak g$ be a differential graded Lie
ideal: it is graded, $d\mathfrak r\subseteq\mathfrak r$, and
$[\mathfrak g,\mathfrak r]\subseteq\mathfrak r$.
The quotient $\mathfrak g/\mathfrak r$ inherits a bracket and a
differential. Let $(\mathfrak r)$ denote the two-sided ideal of
$U(\mathfrak g)$ generated by $\mathfrak r$.

\begin{corollary}[Enveloping quotients]\label{sp:quotient}
There are natural differential graded Hopf and graded Hopf
isomorphisms, respectively,
\[
 U(\mathfrak g)/(\mathfrak r)\cong U(\mathfrak g/\mathfrak r),
 \qquad
 H(U(\mathfrak g)/(\mathfrak r))\cong U(H(\mathfrak g/\mathfrak r)).
\]
Their primitive identifications are
\[
 P(U(\mathfrak g)/(\mathfrak r))\cong\mathfrak g/\mathfrak r,
 \qquad
 H(P(U(\mathfrak g)/(\mathfrak r)))
 \cong P(H(U(\mathfrak g)/(\mathfrak r)))
 \cong H(\mathfrak g/\mathfrak r).
\]
\end{corollary}

\begin{proof}
The differential preserves $(\mathfrak r)$ by the Leibniz rule.
Since its generators are primitive,
$\Delta(\mathfrak r)\subseteq(\mathfrak r)\otimes U+
U\otimes(\mathfrak r)$ and its counit vanishes.
A Lie map from $\mathfrak g$ killing $\mathfrak r$ defines both an
algebra map from $U(\mathfrak g)/(\mathfrak r)$ and one from
$U(\mathfrak g/\mathfrak r)$.
Applying this observation to the two canonical generator maps gives
inverse algebra maps.  They commute with differentials and primitive
coproducts on generators, hence everywhere.  They preserve antipodes
by uniqueness.  Apply Theorems~\ref{sp:primitives} and
\ref{sp:cohomology} to $\mathfrak g/\mathfrak r$ for the other claims.
\end{proof}

The cohomology in this corollary is $H(\mathfrak g/\mathfrak r)$.
It cannot in general be replaced by the quotient of $H(\mathfrak g)$
by the image of $H(\mathfrak r)$.
For the abelian pair $du=v$ above, take $\mathfrak r=\kk v$.
Then $H(\mathfrak g)=0$, but $\mathfrak g/\mathfrak r=\kk u$ has zero
differential and nonzero cohomology.  The enveloping quotient is
$\kk[u]$ with zero differential, whereas $H(U(\mathfrak g))=\kk$.

\subsection{The characteristic and the comparison hypotheses}

Division by two resolves odd squares in the ordered basis.
Division by three gives the all-odd overlap identity from graded Jacobi.
Division by $r!$ supplies symmetrization and the cohomology of symmetric
powers.  Division by $r$ excludes higher-length primitives.
Characteristic zero therefore serves separate roles in the proof.

The all-odd overlap can fail in characteristic three if one assumes
only skew symmetry and \eqref{sp:jacobi}.
Over a field of characteristic three, take odd vectors $x,z$, an even
vector $y$, and brackets
\[
 [x,x]=y,\qquad [x,y]=z,\qquad [y,x]=-z,
\]
with all other basis brackets zero.  Skew symmetry holds.
The only potentially nonzero Jacobi expression has three copies of
$x$, and it is a multiple of $3z$, hence zero.
In its associative quotient, however, $y=2x^2$, which forces
$z=xy-yx=0$.  The vector $z$ is killed, so the ordered-basis conclusion
does not follow from those two axioms in characteristic three.

There is a separate primitive obstruction in any positive
characteristic $p$.  For the ordinary abelian Lie algebra $\kk x$,
$U(\kk x)=\kk[x]$ with $x$ even and primitive.  Then
\[
 \Delta(x^p)=x^p\otimes1+1\otimes x^p
\]
because the intermediate binomial coefficients vanish.
The nonzero primitive $x^p$ lies outside $\kk x$.
These calculations explain why the primitive and cohomology
conclusions above retain their characteristic-zero hypothesis.

For a general differential graded coalgebra $C$, the inclusion
$P(C)\to C$ still gives a map $H(P(C))\to P(H(C))$ whenever the
cohomological tensor target has been specified.
The isomorphism proved here uses an actual differential graded
enveloping presentation with its primitive coproduct.
An associative algebra isomorphism alone does not specify that
coproduct.  Consequently a geometric Hall application requires a
compatible differential graded Hopf comparison before these primitive
identifications apply.
